% ruby-overhang.tex -- ルビ>親のオーバーハングを luatexja から実測
% 都(10) に みやこ(15) = ルビ 5pt 長い。隣が仮名なら食い込め、漢字なら食い込めない
% (AH 既定 except-kanji)。box 構造・kern・glue がどう変わるかを見る。
%   docker run --rm -v "$(pwd -W):/work" -w /work/compare texlive/texlive \
%     lualatex -interaction=nonstopmode ruby-overhang.tex

\documentclass[fontsize=10pt]{jlreq}
\usepackage{luatexja-ruby}
\showboxdepth=10
\showboxbreadth=1000

\begin{document}

% (A) 隣が仮名: の 都(みやこ) に — 両隣へ食い込めるはず
\setbox0=\hbox{の\ruby{都}{みやこ}に}
\showbox0

% (B) 隣が漢字: 京 都(みやこ) 府 — except-kanji で食い込めない → アキが要る
\setbox1=\hbox{京\ruby{都}{みやこ}府}
\showbox1

% (C) 単独 (隣なし=行頭行末相当): 都(みやこ) だけ
\setbox2=\hbox{\ruby{都}{みやこ}}
\showbox2

\end{document}
